\documentclass[12pt, twoside]{article}
\documentclass[12pt, twoside]{article}
\usepackage[letterpaper, margin=1in, headsep=0.2in]{geometry}
\setlength{\headheight}{0.6in}
%\usepackage[english]{babel}
\usepackage[utf8]{inputenc}
\usepackage{microtype}
\usepackage{amsmath}
\usepackage{amssymb}
%\usepackage{amsfonts}
\usepackage{siunitx} %units in math. eg 20\milli\meter
\usepackage{yhmath} % for arcs, overparenth command
\usepackage{tikz} %graphics
\usetikzlibrary{quotes, angles}
\usepackage{graphicx} %consider setting \graphicspath{{images/}}
\usepackage{parskip} %no paragraph indent
\usepackage{enumitem}
\usepackage{multicol}
\usepackage{venndiagram}

\usepackage{fancyhdr}
\pagestyle{fancy}
\fancyhf{}
\renewcommand{\headrulewidth}{0pt} % disable the underline of the header
\raggedbottom
\hfuzz=2mm %suppresses overfull box warnings

\usepackage{hyperref}
\usepackage{float}

\title{Algebra 2}
\author{Chris Huson}
\date{June 2024}

\fancyhead[LE]{\thepage}
\fancyhead[RO]{\thepage \\ Name: \hspace{1.5cm} \,\\}
\fancyhead[LO]{BECA/Huson/Algebra 2: Regents Preparation \\* 20 June 2024}

\begin{document}
\subsubsection*{Prep \#32 - Business and Financial Mathematics}
\begin{enumerate}

\item Julia deposits \$2000 into a savings account that earns 4\% interest per year. The exponential function that models this savings account is $y=2000(1.04)^t$, where $t$ is the time in years. Which equation correctly represents the amount of money in her savings account in terms of the monthly growth rate? 
    %Use the rules of exponents to explain why function that correctly represents the amount of money in her savings account in terms of the \emph{monthly} growth rate is $y=2000(1.0032737)^{12t}$. 
    \begin{multicols}{2}
    \begin{enumerate}
        \item $y=166.67(1.04)^{0.12t}$
        \item $y=2000(1.01)^{t}$
        \item $y=2000(1.0032737)^{12t}$
        \item $y=166.67(1.0032737)^{12t}$
    \end{enumerate}
    \end{multicols} \vspace{0.5cm} %Alg2 Regents Jan2019

\item Jim is looking to buy a vacation home for \$172,600 near his favorite southern beach. The formula to compute a mortgage payment, $M$, is $\displaystyle M=P \cdot \frac{r(1+r)^N}{(1+r)^N-1}$ where $P$ is the principal amount of the loan, $r$ is the monthly interest rate, and $N$ is the number of monthly payments. Jim’s bank offers a monthly interest rate of 0.305\% for a 15-year mortgage.\\*[5pt]
With no down payment, determine Jim’s mortgage payment, rounded to the nearest dollar.\\*[0.5cm]
Algebraically determine and state the down payment, rounded to the \emph{nearest dollar}, that Jim needs to make in order for his mortgage payment to be \$1100. \vspace{0.5cm}
 %Alg2 Regents Jun2017 multiple choice

\item An equation to represent the value of a car after $t$ months of ownership is $\displaystyle v=32,000(0.81)^{\frac{t}{12}}$. Which statement is \emph{not} correct?
\begin{enumerate}
    \item The car lost approximately 19\% of its value each month.
    \item The car maintained approximately 98\% of its value each month.
    \item The value of the car when it was purchased was \$32,000.
    \item The value of the car 1 year after it was purchased was \$25,920.
\end{enumerate} \vspace{0.5cm} %Alg2 Regents Aug2016 MC
        
\item The function below models the average price of gas in a small town since January 1st.
    \[G(t)=-0.0049t^4 + 0.0923t^3 - 0.56t^2 +1.166t+3.23 \text{, where } 0 \leq t \leq 10.\]
    If $G(t)$ is the average price of gas in dollars and $t$ represents the number of months since January 1st, the absolute maximum $G(t)$ reaches over the given domain is about what value, to the nearest cent? \\[1cm]%Alg2 Regents Jan2018

\newpage
 \item The value of a certain small passenger car based on its use in years is modeled by $V(t) =28482.698(0.684)^t$, where $V(t)$ is the value in dollars and $t$ is the time in years. Zach had to take out a loan to purchase the small passenger car. The function $Z(t)=22151.327(0.778)^t$, where $Z(t)$ is measured in dollars, and $t$ is the time in years, models the unpaid amount of Zach's loan over time.\\*[10pt]
Graph $V(t)$ and $Z(t)$ over the interval $0 \leq t \leq 5$, on the set of axes below.
\begin{center}
    \begin{tikzpicture}
    \draw[step=0.25in,gray,very thin] (0,0) grid (12.7,12.7);
    \draw[thick,->] (0,0) -- (13,0);% node[anchor=north west] {x};
    \draw[thick,->] (0,0) -- (0,13);% node[anchor=south east] {y};
    \end{tikzpicture}
\end{center}
State when $V(t)=Z(t)$, to the \emph{nearest hundredth}, and interpret its meaning in the context of the problem.\\*[10pt]
%Alg2 Regents Aug2017 #37


\end{enumerate}
\end{document}